\begin{frame}{자주 묻는 질문 (2/2)}
  \begin{itemize}
    \item \kq{실전에서는 $n$을 얼마로?} --- 후보 4$\sim$10개를
      \textbf{검증 성능}으로 비교. 이 절의 실습이 그 튜닝의 축소판.
      DQN 구현 중엔 버퍼에서 미니배치를 이어붙여 n-step 목표값을
      만드는 변형도 (Week 9 실습에서 마주침).
    \item \kq{``$n$이 크면 분산 큼''은 갱신당 변동이 크다는 뜻인가?}
      --- 그렇다. $n$ 스텝치 전이·보상 무작위성이 합쳐짐(보상 0의
      스텝도 ``어디로 갔는가''의 전이 무작위성은 누적). 같은 총
      정보를 작은 덩어리(잦은 갱신) vs 큰 덩어리(드문 큰 갱신)로
      나누는 것 = ML1 Week 9의 \textbf{배치 크기 트레이드오프}와
      같은 구조.
    \item \kq{적격흔적(7.2)은 정말 같은 것을 계산하나?} --- ``같은
      효과(모든 $n$ 가중합)를 \textbf{다른 순서}로 만든다''가 정확.
  \end{itemize}
\end{frame}
